\documentclass[11pt]{article}
\usepackage{notesstyle}

\begin{document}
\notesheader{3.1}{Introducing Pandas Objects}{Series, DataFrame, and Index --- the three data structures everything else is built on.}

\section{Where Pandas comes from, and why you need it}

\begin{background}
If you have spent your first year of Python mostly with lists, dictionaries, and maybe a little NumPy, here is the mental gap that Pandas fills.

\textbf{NumPy} gives you the \code{ndarray}: a fast, homogeneous, $N$-dimensional grid of numbers. It is brilliant for math, but it has two blind spots that show up the moment you touch \emph{real} data:
\begin{itemize}
  \item \textbf{Positions, not labels.} A NumPy array is addressed by integer position (\code{a[0]}, \code{a[3]}). Real datasets are addressed by \emph{meaning}: the row for ``California,'' the column named ``population.''
  \item \textbf{One dtype, no gaps.} A NumPy array is a single data type with no natural notion of ``missing.'' Real tables mix strings, floats, dates, and holes.
\end{itemize}
\textbf{Pandas} is a thin, powerful layer \emph{on top of} NumPy that attaches \emph{labels} to the axes of an array and gracefully handles missing values. Underneath, your numbers are still NumPy arrays (so you keep the speed); on top, you get a spreadsheet-like, database-like interface.
\end{background}

The convention everyone uses, and the one we use for the rest of these notes:

\begin{lstlisting}
import numpy as np
import pandas as pd
\end{lstlisting}

\begin{intuition}
Think of the three core objects as a ladder of dimensionality, each one a labeled version of a NumPy array:
\begin{itemize}
  \item \textbf{\code{Series}} $=$ a 1-D NumPy array $+$ an explicit index (labels for each element). Like a cross between a NumPy array and a Python dict.
  \item \textbf{\code{DataFrame}} $=$ a 2-D table $=$ an ordered collection of \code{Series} that all \emph{share the same row index}. Like a dict of columns.
  \item \textbf{\code{Index}} $=$ the immutable label array that both of the above use to align data.
\end{itemize}
\end{intuition}

\begin{center}
\begin{tikzpicture}[font=\sffamily\small,
   lbl/.style={font=\sffamily\footnotesize\color{accent}},
   cell/.style={draw=rulegray, minimum width=1.5cm, minimum height=0.7cm},
   idx/.style={cell, fill=bgblue, font=\sffamily\small\color{accent}},
   val/.style={cell, fill=white}]
  % Series
  \node[lbl] at (0,2.2) {\textbf{Series} (1-D + index)};
  \foreach \y/\lab/\v [count=\i] in {3/a/0.25, 2/b/0.50, 1/c/0.75, 0/d/1.00} {
    \node[idx] at (0,\y*0.75) {\lab};
    \node[val] at (1.5,\y*0.75) {\v};
  }
  \node[lbl] at (0,-0.35) {index};
  \node[lbl] at (1.5,-0.35) {values};

  % DataFrame
  \node[lbl] at (6.4,2.2) {\textbf{DataFrame} (aligned columns)};
  \foreach \x/\col [count=\j] in {0/pop, 1/area} {
    \node[idx, fill=bggreen, text=accent2] at (5.6+\x*1.6, 2.25) {\col};
  }
  \foreach \y/\lab [count=\i] in {3/CA, 2/TX, 1/NY, 0/FL} {
    \node[idx] at (4.0,\y*0.75) {\lab};
    \node[val] at (5.6,\y*0.75) {$\bullet$};
    \node[val] at (7.2,\y*0.75) {$\bullet$};
  }
  \node[lbl] at (4.0,-0.35) {shared index};
  \draw[decorate,decoration={brace,amplitude=5pt},accent2] (4.8,2.55) -- (8.0,2.55);
  \node[lbl, text=accent2] at (6.4,3.0) {columns};
\end{tikzpicture}
\end{center}

\section{The Series object}

A \code{Series} is a one-dimensional array of \emph{indexed} data. The simplest way to build one is from a list or NumPy array:

\begin{lstlisting}
data = pd.Series([0.25, 0.5, 0.75, 1.0])
print(data)
\end{lstlisting}

\begin{outblock}
0    0.25
1    0.50
2    0.75
3    1.00
dtype: float64
\end{outblock}

Notice that the printout has \emph{two} columns: the index on the left (here the default integers \code{0..3}) and the values on the right. Those are exactly the two attributes you can pull out:

\begin{lstlisting}
data.values   # -> array([0.25, 0.5 , 0.75, 1.  ])  (a NumPy array)
data.index    # -> RangeIndex(start=0, stop=4, step=1)
\end{lstlisting}

\begin{keyidea}
The \code{values} are a real NumPy array --- that is where the speed lives. The \code{index} is a separate, \emph{explicit} object. Standard NumPy arrays have only an \emph{implicit} integer index; the whole point of a \code{Series} is that you can make the index whatever you want.
\end{keyidea}

\subsection{An index of your own choosing}

Unlike a NumPy array, the index need not be integers, and need not start at zero:

\begin{lstlisting}
data = pd.Series([0.25, 0.5, 0.75, 1.0],
                 index=['a', 'b', 'c', 'd'])
data['b']     # -> 0.5
\end{lstlisting}

This is what people mean when they say a \code{Series} is ``like a dictionary'': you look values up by a label key. The difference from a plain Python \code{dict} is that the keys (the index) are a \emph{typed, ordered} array, which makes slicing and vectorized math possible.

\subsection{Series from a dictionary}

Because a \code{Series} maps keys to values, building one directly from a \code{dict} is natural. Pandas sorts the keys into the index:

\begin{lstlisting}
population_dict = {'California': 39_512_223,
                  'Texas':      28_995_881,
                  'New York':   19_453_561,
                  'Florida':    21_477_737}
population = pd.Series(population_dict)
\end{lstlisting}

\begin{outbox}
California \quad 39512223 \\
Florida \quad\;\; 21477737 \\
New York \quad\, 19453561 \\
Texas \quad\quad 28995881 \\
dtype: int64
\end{outbox}

Now it behaves like a dictionary --- \code{population['California']} returns the value --- \emph{and} like an array, supporting slicing such as \code{population['California':'New York']}.

\begin{pitfall}
\textbf{Label slicing includes the right endpoint.} With the default integer positions, \code{data[0:2]} gives 2 elements (the usual Python ``stop is exclusive'' rule). But when you slice by \emph{explicit labels}, e.g. \code{data['a':'c']}, Pandas includes \code{'c'}. Two different conventions live in the same object --- this is exactly the ambiguity that \code{loc} and \code{iloc} (next section) exist to resolve.
\end{pitfall}

\section{The DataFrame object}

If a \code{Series} is a labeled 1-D array, a \code{DataFrame} is a labeled 2-D array: a table whose rows share a common index and whose columns each have a name. Equivalently, and more usefully: \emph{a \code{DataFrame} is a dictionary of \code{Series} that all share the same index.}

Suppose we have a second \code{Series} for area, indexed by the same states:

\begin{lstlisting}
area = pd.Series({'California': 423_967, 'Texas': 695_662,
                  'New York': 141_297, 'Florida': 170_312})
states = pd.DataFrame({'population': population, 'area': area})
\end{lstlisting}

\begin{center}
\renewcommand{\arraystretch}{1.2}
\begin{tabular}{l r r}
\toprule
 & \textbf{population} & \textbf{area} \\
\midrule
California & 39512223 & 423967 \\
Florida    & 21477737 & 170312 \\
New York   & 19453561 & 141297 \\
Texas      & 28995881 & 695662 \\
\bottomrule
\end{tabular}
\end{center}

Like the \code{Series}, the \code{DataFrame} exposes its structure through attributes:

\begin{lstlisting}
states.index     # the row labels: ['California','Florida','New York','Texas']
states.columns   # the column labels: ['population', 'area']
states.values    # the underlying 2-D NumPy array
\end{lstlisting}

\begin{intuition}
There are two ways to read a \code{DataFrame}, and you will switch between them constantly:
\begin{itemize}
  \item \textbf{As a dict of columns:} \code{states['area']} returns the \code{area} \code{Series}. This is the \emph{default} --- a single key indexes a \emph{column}, not a row. (This trips up people coming from NumPy, where \code{a[0]} is a row.)
  \item \textbf{As an enhanced 2-D array:} you can transpose it (\code{states.T}), do matrix-style math, and use the \code{loc}/\code{iloc} accessors for row/column selection.
\end{itemize}
\end{intuition}

\subsection{Many ways to construct a DataFrame}

You will meet all of these in the wild:

\begin{lstlisting}
# 1. From a single Series (one column)
pd.DataFrame(population, columns=['population'])

# 2. From a list of dicts (each dict is a row)
pd.DataFrame([{'a': 1, 'b': 2}, {'a': 3, 'b': 4}])

# 3. From a dict of Series (each value is a column)  <- most common
pd.DataFrame({'population': population, 'area': area})

# 4. From a 2-D NumPy array, naming axes explicitly
pd.DataFrame(np.random.rand(3, 2),
             columns=['foo', 'bar'],
             index=['x', 'y', 'z'])
\end{lstlisting}

\begin{pitfall}
When you build from a \emph{list of dicts} and the dicts have different keys, Pandas fills the missing entries with \code{NaN} rather than erroring. That is convenient but can silently hide a typo in a key name --- always check \code{df.columns} after construction.
\end{pitfall}

\section{The Index object}

Both \code{Series} and \code{DataFrame} contain an \code{Index}: a curious little object that is best thought of in two ways at once.

\subsection{Index as an immutable array}

You can slice and index it like a NumPy array, but you \emph{cannot} mutate it:

\begin{lstlisting}
ind = pd.Index([2, 3, 5, 7, 11])
ind[1]      # -> 3
ind[::2]    # -> Int64Index([2, 5, 11])
ind[1] = 0  # TypeError: Index does not support mutable operations
\end{lstlisting}

\begin{background}
\textbf{Why on earth make the index immutable?} Because several \code{Series}/\code{DataFrame} objects can \emph{share} the same index, and many operations \emph{align} on it. If you could mutate an index in place, you could silently corrupt every object that shares it --- like changing a primary key in a database while other tables still reference the old value. Immutability makes sharing safe and makes alignment (the topic of section 3.3) predictable.
\end{background}

\subsection{Index as an ordered set}

Pandas indices support set operations, which is what makes joining and aligning datasets possible:

\begin{lstlisting}
indA = pd.Index([1, 3, 5, 7, 9])
indB = pd.Index([2, 3, 5, 7, 11])
indA.intersection(indB)   # -> [3, 5, 7]
indA.union(indB)          # -> [1, 2, 3, 5, 7, 9, 11]
indA.symmetric_difference(indB)  # -> [1, 2, 9, 11]
\end{lstlisting}

\begin{keyidea}
Keep this picture in your head for the rest of Chapter 3: \textbf{every Pandas operation that combines data --- arithmetic, merging, grouping --- works by lining up \emph{labels} via these \code{Index} objects, not by lining up \emph{positions}.} Once labels drive everything, ``missing on one side'' becomes a normal, expected case rather than an error, which is why missing-data handling (section 3.4) is woven so deeply into the library.
\end{keyidea}

\section*{Section recap}
\addcontentsline{toc}{section}{Section recap}
\begin{itemize}
  \item Pandas $=$ NumPy speed $+$ labeled axes $+$ missing-data support.
  \item \code{Series}: 1-D values $+$ an explicit \code{Index}; behaves like a dict \emph{and} an array.
  \item \code{DataFrame}: aligned columns sharing one row index; behaves like a dict of \code{Series}. A bare key selects a \emph{column}.
  \item \code{Index}: immutable, set-like; it is the machinery that aligns data across objects.
  \item Watch the slicing rule: label slices include the endpoint, positional slices do not.
\end{itemize}

\end{document}
